%! Exponential and Logarithmic Equations and Inequalities — Unit 5, Lesson 5.6 — Lesson Plan
\documentclass[10pt]{article}
\usepackage{precalculus-article}
\usepackage{precalculus-boxes}
\usepackage{graphicx}
\graphicspath{{images/}}
\newcommand{\TallMath}[1]{$\displaystyle #1\rule[-1.4em]{0pt}{3.2em}$}
\newcommand{\CourseName}{Precalculus}
\newcommand{\SchoolYear}{2026--2027}
\newcommand{\MeetingLength}{55 minutes}
\newcommand{\UnitNumberName}{Unit 5: Logarithmic Functions \quad}
\newcommand{\LessonNumberName}{Lesson 5.6: Exponential and Logarithmic Equations and Inequalities}

\begin{document}
\begin{center}
    {\Large\bfseries \CourseName: \SchoolYear \\
    {\small\bfseries \UnitNumberName \LessonNumberName}}
\end{center}

\begin{tcolorbox}[colback=sky, colframe=navy, boxrule=0.9pt, arc=2mm,
    left=2mm, right=2mm, top=1.5mm, bottom=1.5mm]
    \textbf{Primary Objective:} Students solve exponential and logarithmic equations and
    inequalities by applying properties of exponents and logarithms and the inverse
    relationship between exponential and logarithmic functions; check results for extraneous
    solutions; use the identity $b^x = e^{(\ln b)\cdot x}$ to rewrite exponential expressions;
    and construct inverse functions for transformed exponential and logarithmic functions.
    \hfill\textit{\small Skills: AP 1.A, AP 1.B, AP 1.C \quad LO 2.13.A, LO 2.13.B}
\end{tcolorbox}

\renewcommand{\arraystretch}{1.6}
\begin{skillbox}[Priority Ideas \& Skills]{goldbox}
\begin{tabularx}{\linewidth}{@{}X X@{}}
\textbf{AP Skills} &
\textbf{Key Understandings} \\
\midrule
\textbf{AP 1.A} --- Solve equations and inequalities. &
\begin{itemize}[leftmargin=1.2em,itemsep=2pt,topsep=0pt]
  \item Properties of exponents and logarithms and the inverse relationship can be used to
        solve exponential and logarithmic equations and inequalities. (EK 2.13.A.1)
  \item Results must be examined for extraneous solutions. (EK 2.13.A.2)
\end{itemize} \\
\textbf{AP 1.B} --- Construct equivalent representations. &
\begin{itemize}[leftmargin=1.2em,itemsep=2pt,topsep=0pt]
  \item $b^x = e^{(\ln b)\cdot x}$ rewrites any exponential in natural-base form.
        (EK 2.13.A.3)
\end{itemize} \\
\textbf{AP 1.C} --- Construct new functions using inverses. &
\begin{itemize}[leftmargin=1.2em,itemsep=2pt,topsep=0pt]
  \item Inverse of $f(x)=ab^{x-h}+k$: reverse the operation sequence. (EK 2.13.B.1)
  \item Inverse of $f(x)=a\log_b(x-h)+k$: reverse the operation sequence. (EK 2.13.B.2)
\end{itemize} \\
\end{tabularx}
\end{skillbox}

\begin{skillbox}[{Vocabulary, Concepts, \& Theorems}]{greenbox}
\renewcommand{\arraystretch}{1.4}
\begin{tabularx}{\linewidth}{@{} >{\bfseries\color{navy}}p{0.28\linewidth} X @{}}
extraneous solution & A value obtained algebraically that does not satisfy the original
                     equation because the argument of a logarithm must be positive (or other
                     domain restriction). Always substitute back to verify. \\
one-to-one property & If $b^a = b^c$ (same base, $b>0$, $b\ne1$), then $a=c$.
                     If $\log_b a = \log_b c$, then $a=c$. \\
natural base identity & $b^x = e^{(\ln b)\cdot x}$ for any $b>0$, $b\ne1$. Unifies all
                     exponential functions under the natural base. (EK 2.13.A.3) \\
inverse operations & To find $f^{-1}$, swap $x$ and $y$, then undo each operation of $f$
                    in reverse order. \\
\end{tabularx}
\end{skillbox}

\begin{skillbox}[Activate Prior Knowledge \& Spiral Review]{sky}
    Spiral review (warm-up) covers: solving $2^x = 32$ and $3^x = \tfrac{1}{27}$ using the
    one-to-one property; evaluating $g(f(8))$ where $f(x)=\log_2 x$ and $g(x)=2^x$ to
    reinforce inverse cancellation (answer: 8); and identifying the two transformations applied
    to the basic exponential $2^x$ in $f(x)=5\cdot 2^x+1$ (vertical stretch by 5, then
    vertical shift up 1). These connect to EK 2.13.A.1 (one-to-one property),
    EK 2.13.A.2 (inverse cancellation confirms correct solutions), and EK 2.13.B.1
    (recognizing transformation structure before inverting).
\end{skillbox}

\begin{skillbox}[Hook]{sky}
    Write on the board: \textbf{A doctor prescribes a drug with a half-life of 4 hours.
    The initial dose is 200 mg. After how many hours will less than 10 mg remain?}

    Ask students to write the model: $D(t) = 200\cdot\bigl(\tfrac{1}{2}\bigr)^{t/4}$.
    Then pose: ``How do we solve $200\cdot\bigl(\tfrac{1}{2}\bigr)^{t/4} < 10$ for $t$?''
    Students quickly see that guess-and-check is tedious. Ask: ``What tool lets us isolate
    an exponent?'' This motivates taking logarithms of both sides (EK 2.13.A.1).
\end{skillbox}

\begin{skillbox}[Lesson]{sky}
\begin{multicols}{2}
\textbf{Part 1 (10 min): Solving Exponential Equations (EK 2.13.A.1).}

\textit{Method 1 --- same-base / one-to-one property:} $2^x = 2^5 \Rightarrow x = 5$.

\textit{Method 2 --- take log of both sides:} $3^x = 15 \Rightarrow \log(3^x) = \log 15
\Rightarrow x\log 3 = \log 15 \Rightarrow x = \tfrac{\log 15}{\log 3} \approx 2.46$.

Note: any log base works; the ratio $\tfrac{\log c}{\log b} = \log_b c$.

Work through $5^x = 100$ with students filling guided notes.

\columnbreak

\textbf{Part 2 (10 min): Solving Logarithmic Equations and Checking for Extraneous
Solutions (EK 2.13.A.1--A.2).}

$\log_2 x = 4 \Rightarrow 2^4 = 16 \Rightarrow x = 16$. (Check: $\log_2 16 = 4$. \checkmark)

$\log_5 x + \log_5(x-4) = 1$: combine logs $\Rightarrow \log_5(x(x-4))=1 \Rightarrow
x(x-4)=5 \Rightarrow x^2-4x-5=0 \Rightarrow (x-5)(x+1)=0$.
Check: $x=-1$ makes $\log_5(-1)$ undefined --- extraneous. Answer: $x=5$ only.
\end{multicols}
\end{skillbox}

\begin{skillbox}[Lesson (cont.)]{sky}
\begin{multicols}{2}
\textbf{Part 3 (8 min): Natural Base Identity (EK 2.13.A.3).}

$b^x = e^{(\ln b)\cdot x}$. Example: $2^x = e^{(\ln 2)\cdot x} \approx e^{0.693x}$.
Discuss how this unifies exponential functions. Fill in: $3^x = e^{(\blank{1.2cm})\cdot x}$
and $5^x = e^{(\blank{1.2cm})\cdot x}$.

\columnbreak

\textbf{Part 4 (12 min): Inverses of Transformed Functions (EK 2.13.B.1--B.2).}

\textit{Exponential:} $f(x)=3\cdot2^{x-1}+4$. Swap, isolate step by step: subtract 4,
divide by 3, apply $\log_2$, add 1. $f^{-1}(x)=\log_2\!\tfrac{x-4}{3}+1$.

\textit{Logarithmic:} $f(x)=2\log_3(x+1)-5$. Swap, isolate: add 5, divide by 2,
exponentiate base 3, subtract 1. $f^{-1}(x)=3^{(x+5)/2}-1$.

For each, state domain and range of $f$ and $f^{-1}$.
\end{multicols}
\end{skillbox}

\begin{skillbox}[Active Monitoring]{redbox}
\begin{itemize}[leftmargin=1.2em,itemsep=2pt,topsep=0pt]
  \item After Part 1: cold-call ``Why does any log base work when taking log of both sides?''
        Expect: the ratio $\log c/\log b$ is independent of the base chosen.
  \item Watch for students writing $x\log 3 = \log 15$ but then computing $x = \log 15 - \log 3$
        (subtracting instead of dividing) --- address the power rule of logs immediately.
  \item After Part 2: ask ``How do you know $x=-1$ is extraneous?'' Expect: the argument
        $\log_5(-1)$ is undefined because log domain requires positive input.
  \item During Part 4: confirm students start by swapping $x \leftrightarrow y$, not by
        ``plugging in'' the inverse formula directly. Both methods should give the same result.
\end{itemize}
\end{skillbox}

\begin{skillbox}[Group Work \& Differentiation]{redbox}
\begin{itemize}[leftmargin=1.2em,itemsep=2pt,topsep=0pt]
  \item \textbf{Tier R (Remediate):} Solve four basic exponential and log equations with
        integer solutions; verify each solution in the original equation; identify domain
        restriction for log equations.
  \item \textbf{Tier A (Apply):} Solve the half-life inequality from the hook; solve a
        combined-log equation checking for extraneous solutions; apply the natural base
        identity.
  \item \textbf{Tier E (Extension):} Solve an exponential inequality with competing bases;
        find inverses of a transformed exponential and a transformed logarithmic function,
        each with full domain/range analysis.
\end{itemize}
\end{skillbox}

\begin{skillbox}[Individual Work \& Assessment]{redbox}
Exit ticket (2 questions, 1 page):
\begin{enumerate}[leftmargin=1.5em,itemsep=2pt,topsep=2pt]
  \item Solve $5^x = 80$; give exact form and decimal approximation to 3 decimal places.
        (Tests EK 2.13.A.1 --- taking log of both sides.)
  \item Solve $\log_4(x+2)=3$; check for domain restrictions and state whether the
        solution is valid. (Tests EK 2.13.A.1--A.2 --- convert and check extraneous.)
\end{enumerate}

Sort by result: students missing \#1 need more practice with the log-of-both-sides
technique; students missing \#2 need the habit of substituting back to check domain.
\end{skillbox}

\begin{skillbox}[Reinforcement \& Extension]{goldbox}
\textbf{Homework (6 problems + extension):} Solve exponential equations using same-base
and log methods; solve log equations with extraneous-solution checks; population growth
model requiring logarithms; natural base identity rewrites; find inverse of a transformed
exponential; AP-style MC using the product rule for logarithms; extension inequality.

\textbf{Preview of Lesson 2.14:} Connecting all representations of exponential and
logarithmic functions --- graphs, tables, equations, and contextual models --- in preparation
for AP exam multi-part problems.
\end{skillbox}

\end{document}
